\documentclass[11pt]{article}
\usepackage{ap-physc}
\course{AP Physics C: Electricity and Magnetism}
\headcourse{AP Physics C: E\&M}
\headpart{Unit 12 \textperiodcentered\ Study Guide}
\title{Unit 12: Magnetic Fields and Electromagnetism}
\author{Paul Wang}
\date{November 2026}
\docpart{Unit 12 \textperiodcentered\ Study Guide \textperiodcentered\ 10--20\% of MCQ}

\begin{document}
\maketitle

\begin{abstract}
Unit 12 covers magnetism produced by and acting on moving charges. It covers magnetic fields and Gauss's law for magnetism (12.1); the force $q\vec v\times\vec B$ on moving charges and $I\vec L\times\vec B$ on currents, circular motion, velocity selectors, and torque on current loops (12.2); the Biot--Savart law for wires, loops, and arcs (12.3); and Ampère's law for wires, solenoids, toroids, and coaxial cables (12.4). Each section ends with check questions; the key is on the last page.
\end{abstract}

\scoringline{Check questions: \blank\ / \qtotal \qquad Target: 5+ \qquad Weakest topic: \blank[8em]}

\section{Magnetic fields (12.1)}\label{sec:1}

Magnetic fields are produced by moving charges: currents, and the spins and orbits of electrons in magnetic materials. Field lines leave a magnet's north pole and enter its south pole, but unlike electric field lines they never start or end: they form closed loops. There are no magnetic monopoles, so
\[ \oint\vec B\cdot d\vec A = 0 \quad \text{(Gauss's law for magnetism)}. \]
The SI unit is the tesla, $1\ \mathrm{T} = 1\ \mathrm{N/(A\cdot m)}$. Earth's field is about $5\times10^{-5}$ T; a classroom bar magnet, about $10^{-2}$ T near its poles.
\begin{trap}
Cutting a bar magnet in half gives two smaller magnets, each with both poles, never an isolated pole.
\end{trap}

\subsubsection{Check}\mbox{}

\begin{q}{2}{12.1 Magnetic flux through a closed surface}{Magnetic field lines form closed loops; every line that leaves the surface also enters it somewhere (inside the magnet the lines run from S to N). Net flux is always zero.}
A closed surface encloses only the north-pole end of a bar magnet. What is the net magnetic flux through the surface?
\ch{Positive, because field lines leave the north pole}{Zero}{Negative, because field lines enter the south pole}{It depends on the magnet's strength}
\end{q}

\section{Magnetic forces on moving charges and currents (12.2)}\label{sec:2}

\[ \vec F = q\vec v\times\vec B, \qquad F = |q|vB\sin\theta, \qquad \vec F = I\vec L\times\vec B\ \ \text{(straight wire)}. \]
\begin{itemize}
\item The force is perpendicular to both $\vec v$ and $\vec B$ (right-hand rule; flip it for negative charges). Because $\vec F \perp \vec v$, a magnetic force does \textbf{no work} and never changes a particle's speed.
\item A charge moving perpendicular to a uniform $\vec B$ moves in a circle: $qvB = mv^2/r$, so $r = \dfrac{mv}{qB}$, and the period $T = \dfrac{2\pi m}{qB}$ does not depend on speed.
\item Crossed $\vec E$ and $\vec B$ fields pass undeflected only particles with $v = E/B$ (velocity selector).
\item A current loop in a uniform field feels no net force but a torque $\vec\tau = \vec\mu\times\vec B$, with magnetic moment $\mu = NIA$ along the loop's normal (right-hand rule with the current).
\end{itemize}

\subsubsection{Check}\mbox{}

\begin{q}{4}{12.2 Direction of force}{$\vec F = q\vec v\times\vec B \propto \hat\imath\times\hat\jmath = +\hat k$; the proton is positive, so the force is $+z$.}
A proton moves in the $+x$ direction through a magnetic field that points in the $+y$ direction. Which way is the magnetic force on it?
\chh{$-z$}{$+y$}{$-x$}{$+z$}
\end{q}
\begin{q}{1}{12.2 Work by a magnetic force}{The magnetic force is always perpendicular to the velocity, so it does no work: it changes direction, not speed.}
A charged particle moves through a region with only a uniform magnetic field. What happens to its speed?
\ch{It stays constant.}{It increases.}{It decreases.}{It increases if the charge is positive and decreases if negative.}
\end{q}

\section{Fields of currents and the Biot--Savart law (12.3)}\label{sec:3}

\[ d\vec B = \frac{\mu_0}{4\pi}\frac{I\,d\vec l\times\hat r}{r^2}, \qquad \mu_0 = 4\pi\times10^{-7}\ \mathrm{T\,m/A}. \]
\begin{compare}{Source}{Field}{Where}
Long straight wire & $\dfrac{\mu_0 I}{2\pi r}$, circling the wire & distance $r$ \\
Circular loop, radius $R$ & $\dfrac{\mu_0 I}{2R}$, along the axis & at the center \\
Arc of angle $\theta$ (rad) & $\dfrac{\mu_0 I\theta}{4\pi R}$ & at the center of the arc \\
Straight segment pointing at the point & $0$ & any point on its line ($d\vec l \parallel \hat r$) \\
\end{compare}
\textbf{Parallel wires} exert forces per length $\dfrac{F}{L} = \dfrac{\mu_0 I_1I_2}{2\pi d}$: currents in the same direction attract, opposite directions repel (the reverse of charges).
\begin{trap}
Use the right-hand grip rule for direction: thumb along the current, fingers curl the way $\vec B$ circles the wire.
\end{trap}

\subsubsection{Check}\mbox{}

\begin{q}{3}{12.3 Long straight wire}{$B = \dfrac{\mu_0 I}{2\pi r} = \dfrac{(2\times10^{-7})(10)}{0.020} = 1.0\times10^{-4}$ T.}
A long straight wire carries 10 A. What is the magnetic field 2.0 cm from it?
\chh{$5.0\times10^{-3}$ T}{$3.1\times10^{-4}$ T}{$1.0\times10^{-4}$ T}{$1.0\times10^{-6}$ T}
\end{q}
\begin{q}{3}{12.3 Parallel currents}{Opposite currents repel; same-direction currents attract. Check with the right-hand rule: the field of one wire, crossed with the other's current, points away.}
Two long parallel wires carry currents in opposite directions. The wires
\ch{attract each other}{exert no force on each other}{repel each other}{twist to become perpendicular}
\end{q}

\section{Amp\`ere's law (12.4)}\label{sec:4}

\[ \oint\vec B\cdot d\vec l = \mu_0 I_{enc}. \]
Like Gauss's law, it always holds but gives $B$ only with symmetry: choose an Amp\`erian loop along which $B$ is constant and parallel to the path (or perpendicular, contributing nothing).
\begin{compare}{Current distribution}{Amp\`erian loop}{Result}
Long wire (outside) & circle of radius $r$ & $B = \mu_0 I/(2\pi r)$ \\
Inside a long wire, uniform current, radius $R$ & circle $r < R$; $I_{enc} = Ir^2/R^2$ & $B = \dfrac{\mu_0 Ir}{2\pi R^2}$ \\
Long solenoid, $n$ turns per length & rectangle, one side inside & $B = \mu_0 nI$ inside, $\approx 0$ outside \\
Toroid, $N$ turns & circle inside the core & $B = \dfrac{\mu_0 NI}{2\pi r}$ \\
\end{compare}
A coaxial cable with equal and opposite currents has $I_{enc} = 0$ outside it, so no field leaks out.

\subsubsection{Check}\mbox{}

\begin{q}{1}{12.4 Solenoid}{$B = \mu_0 nI \propto n$.}
A long solenoid's number of turns per meter is doubled while the current stays the same. What happens to the field inside?
\chh{It doubles.}{It halves.}{It quadruples.}{It is unchanged.}
\end{q}

\makekey
\end{document}
